% Part: history 
% Chapter: biographies 
% Section: julia-robinson
\documentclass[../../../include/open-logic-section]{subfiles}

\begin{document}

\olfileid{his}{bio}{rob}

\olsection{جوليا رابنسن}

\olphoto{robinson-julia}{جوليا رابنسن}

جوليا بومن رابنسن يوه امريکايۍ رياضيپوهه
وه۔ هغه په عمده توګه د پرېکړې پر مسئلو
د خپل کار، او تر ټولو زيات د هېلبرټ
د لسمې مسئلې په حل کښې د خپلو مرستو
له امله پېژندل کېږي۔ رابنسن د 1919
د دسمبر پر 8مه د ميزوري په سېنټ
لويس کښې زېږېدلې وه۔ رابنسن يادوي
چې له ماشومتوبه د عددونو سره شوق
لرلو \citep[4]{Reid1986}۔ په نهه
کلنۍ کښې په سرې تبې اخته شوه او
د روماتيکي تبې څو ځله تکراري
حملې ئې تېرې کړې۔ دې هغه اړ کړه
چې ډېر وخت په کټ کښې تېر کړي،
او په زده کړه کښې وروسته شوه۔
که څه هم د شخصي ښوونکو په مرسته
بېرته له نورو سره برابره شوه،
د ناروغۍ بدني اغېزو ئې پر ژوند
تلپاتې اثر وکړ۔

د ماشومتوب له ستونزو سره سره رابنسن
له ثانوي ښوونځي څخه په رياضياتو او
ساينس کښې له څو انعامونو سره فارغه
شوه۔ پوهنتوني زده کړې ئې د سېن
ډياګو په دولتي کالج کښې پيل کړې،
او په وروستي کال کښې د کاليفورنيا
برکلي پوهنتون ته لاړه۔ هلته د
رياضيپوه رافايل رابنسن تر اغېز
لاندې راغله۔ ښه ملګري شول او په
1941 کښې ئې واده وکړ۔ د يو استاد
د مېرمنې په توګه رابنسن په برکلي
کښې د رياضياتو په څانګه کښې له
تدريسه منع وه۔ که څه هم د
رياضياتو ټولګي ئې لا هم اورېدل،
هيله ئې لرله چې پوهنتون پرېږدي
او کورنۍ جوړه کړي۔ خو له واده
څه موده وروسته رابنسن په سينه‌بغل
اخته شوه۔ ورته وويل شول چې د
ماشومتوب د روماتيکي تبې له امله
ئې پر زړه ډېر داغي نسج راټول
شوے دے۔ د دې نسج د سختوالي له
امله ډاکټر وړاندوينه وکړه چې
له څلوېښتو کلونو زيات ژوند به
نۀ کوي، او ورته مشوره ورکړے
شوه چې اولاد ونۀ لري
\citep[13]{Reid1986}۔

رابنسن د ډېر وخت دپاره خپه وه، خو
بالاخره ئې پرېکړه وکړه چې د
رياضياتو زده کړه جاري وساتي۔ برکلي
ته ستنه شوه او په 1948 کښې ئې
د الفرېډ تارسکي تر لارښوونې
لاندې دوکتورا بشپړه کړه۔ تارسکي
ښودلې وه چې د حقيقي عددونو د
لومړي ترتيب تيوري د پرېکړې وړ
ده، او د ګوډل له کاره دا پايله
راتله چې د طبيعي عددونو د
لومړي ترتيب تيوري د پرېکړې
وړ نۀ ده۔ دا يوه لويه خلاصه
مسئله وه چې د نسبي عددونو
د لومړي ترتيب تيوري د پرېکړې
وړ ده که نۀ۔ رابنسن په خپله
رساله \citeyearpar{Robinson1949}
کښې ثابته کړه چې نۀ ده۔

رابنسن چې د پرېکړې له مسئلو سره ئې
شوق لرلو، بيا ئې د هېلبرټ د لسمې
مسئلې د حل هڅه وکړه۔ دا مسئله
د هغو 23 مشهورو رياضيکي مسئلو
له لړ څخه وه چې ډېوېډ هېلبرټ په
1900 کښې وړاندې کړې وې۔ لسمه
مسئله پوښتي چې آيا داسې الګوريتم
شته چې په متناهي وخت کښې ځواب
ورکړي چې د صحيح عددي ضريبونو
يوه کثيرالحدي معادله، لکه
$3x^2 - 2y +3 = 0$، په صحيح
عددونو کښې حل لري که نۀ۔ داسې
پوښتنې \emph{ديوفانټي مسئلې}
نومېږي۔ له ځينو لومړنيو برياوو
وروسته رابنسن له مارټن ډېوس او
هېلري پټنم سره ګډ کار پيل کړ،
چې هغوی هم پر دې مسئله کار
کاوۀ۔ دوی وښودله چې نمايي
ديوفانټي مسئلې، چې نامعلومونه
پکې د نماوو په توګه هم راتلے
شي، د پرېکړې وړ نۀ دي۔ او
وښودله چې يوه ځانګړې اټکلي
دعوه، چې وروسته «J.R.» ونومول
شوه، دا پايله لري چې د
هېلبرټ لسمه مسئله د پرېکړې
وړ نۀ ده \citep{DavisPutnamRobinson1961}۔
رابنسن په ټولو 1960s کلونو
کښې پر دې مسئله کار جاري
وساتۀ۔ په 1970 کښې ځوان
روسي رياضيپوه يوري ماتياسېوېچ
بالاخره د J.R. فرضيه ثابته
کړه۔ ګډه نتيجه اوس د
ماتياسېوېچ--رابنسن--ډېوس--پټنم
قضيه، يا په لنډه MRDP قضيه،
نومېږي۔ ماتياسېوېچ او رابنسن
ملګري شول او پر څو مقالو ئې
ګډ کار وکړ۔ رابنسن يو ځل
ماتياسېوېچ ته په يوه ليک کښې
وليکل: «اصل کښې زه ډېره خوشحاله
يم چې په ګډ کار، سره له دې
چې زرګونه ميله لرې يو، په
ښکاره تر هغه ډېر پرمختګ کوو
چې زموږ هر يوه يوازې کولے
شو» \citep[45]{Matijasevich1992}۔

رابنسن د امريکا د رياضيکي ټولنې
لومړۍ ښځينه صدره وه، او لومړۍ
ښځه وه چې د علومو ملي اکاډمۍ
ته وټاکل شوه۔ هغه د وينې د
سرطان له تشخيص وروسته د 1985
د جولای پر 30مه د 65 کلونو
په عمر وفات شوه۔
\textbf{د سرچينې د دعوې يادونه (OLSTH-043):}
د اکاډمۍ د ټولو غړو په هکله د
«لومړۍ ښځې» چاپي دعوه ډېره
پراخه ده۔ کره دعوه دا ده چې
رابنسن د همدې اکاډمۍ د رياضياتو
څانګې ته ټاکل شوې لومړۍ ښځه
وه، لکه د لاندې لوست د فېفرمان
په ژوندليکي يادښت کښې چې
بيان شوي دي۔ د اکاډمۍ خپله
تاريخچه فلورنس سابن د ټولې
اکاډمۍ لومړۍ ښځينه غړې بولي۔

\begin{reading}
د رابنسن رياضيکي ليکنې د هغې په
\textit{راټولو شويو ليکنو} کښې
شته \citep{Robinson1996}۔ په دې
کښې د علومو ملي اکاډمۍ ژوندليکي
يادښت \citep{Feferman1994} هم
بيا چاپ شوے دے۔ د رابنسن مشرې
خور کانسټنس ريډ د مرکو پر بنسټ
«د جوليا ځانليک» خپور کړے دے
\citep{Reid1986}، او يو بشپړ
ژوندليکي يادښت هم
\citep{Reid1996}۔ د رابنسن او
د هېلبرټ د لسمې مسئلې په هکله
د يو لنډ مستند فلم لارښوونه
جورج چيچري کړې ده
\citep{Csicsery2016}۔ له رابنسن
سره د يوري ماتياسېوېچ د ګډ
کار، او د هغه پر کار د هغې
د اغېز د لنډ يادښت دپاره
\citep{Matijasevich1992} وګورئ۔
\end{reading}

\end{document}
